% Owner: Member 1 (Task.md, Report Ownership). Text from the SAT draft of 2026-11-14.

\section{Mathematical Formulations \& SAT Encodings}

\subsection{Order Encoding}
Order Encoding utilizes boolean variables $L(u, i) = \text{True} \iff f(u) \le i$. 
\begin{itemize}
    \item \textbf{Boundary Condition:} $L(u, n) = \text{True}, \quad \forall u \in V$
    \item \textbf{Monotonicity Constraints:} $\neg L(u, i) \lor L(u, i+1), \quad \forall i \in \{1, \dots, n-1\}$
    \item \textbf{Bijection / At-Most-One (AMO):} $\neg L(u, i) \lor \neg L(v, i) \lor L(u, i-1) \lor L(v, i-1)$
    \item \textbf{Edge Distance Constraints:} $\neg L(u, k) \lor L(v, k + UB), \quad \forall (u,v) \in E$
\end{itemize}

\subsection{Pairwise Encoding}
Pairwise Encoding defines boolean variables $X(u, i) = \text{True} \iff f(u) = i$.
\begin{itemize}
    \item \textbf{Exactly One Position per Vertex:} Each vertex is assigned exactly one label.
    \item \textbf{Exactly One Vertex per Position:} Each label is assigned to exactly one vertex.
    \item \textbf{Edge Distance Constraints:} $\neg X(u, k) \lor \bigvee_{j = \max(1, k-UB)}^{\min(n, k+UB)} X(v, j), \quad \forall (u,v) \in E$
\end{itemize}

\section{Solving Strategies}

\subsection{Repeated SAT Solver}
At each upper bound evaluation step $UB = k$, a brand-new SAT formula is constructed and passed to the CaDiCaL solver engine from scratch.

\subsection{Incremental SAT Solver}
The solver is initialized once with base structural constraints. For each iteration testing $UB = k$, distance constraints are conditionally activated via helper literals $A_k$ (Activation Literals) in the form:
\begin{equation}
    (\text{Edge Clause}) \lor \neg A_k
\end{equation}
During solving, invoking \texttt{solver.solve(assumptions=[A\_k])} enables the reuse of learned clauses across consecutive bound searches.